\section*{الفصل \ref{ch:g11:infrared} --- مطيافية الأشعة تحت الحمراء}

\begin{solution}{exo:g11:infrared:1}
$\lambda = 1/3300 = \qty{3.03e-4}{cm} = \qty{3.03}{\micro m}$؛
$\lambda = 1/1050 = \qty{9.52e-4}{cm} = \qty{9.52}{\micro m}$. والحزمة عند
\qty{3300}{cm^{-1}}، ذات \omterm{def:g11:infrared:wavenumber}{العدد الموجي} الأكبر، تعود إلى الاهتزاز
الأسرع.
\end{solution}

\begin{solution}{exo:g11:infrared:2}
$\qty{10}{\micro m} = \qty{1.0e-3}{cm}$، إذن $\sigma = 1/\num{1.0e-3} =
\qty{1000}{cm^{-1}}$.
\end{solution}

\begin{solution}{exo:g11:infrared:3}
الرابطة \ce{C=O} \omterm{def:g11:functional-groups:carbonyl}{لكيتون}؛ ورابطة \ce{O-H} \omterm{def:g11:functional-groups:alcohol}{لكحول} مرتبطة بروابط هيدروجينية؛
وروابط \ce{C-H} لسلسلة كربونية.
\end{solution}

\begin{solution}{exo:g11:infrared:4}
الحزمة عند \qty{1715}{cm^{-1}}: الرابطة \ce{C=O}.
\end{solution}

\begin{solution}{exo:g11:infrared:5}
يقيس الجهاز الضوء الذي يجتاز العيّنة؛ فحيث تمتص رابطةٌ يجتاز ضوء أقل
وينخفض المنحنى. و$T =
\qty{100}{\percent}$ تعني أن شيئًا من ضوء هذا \omterm{def:g11:infrared:wavenumber}{العدد الموجي} لم
يُمتص.
\end{solution}

\begin{solution}{exo:g11:infrared:6}
للأول حزمة \ce{C=O} عند \qty{1715}{cm^{-1}} ولا حزمة
\ce{O-H}. فهو \omterm{def:g11:functional-groups:carbonyl}{كيتون}: \ce{CH3-CO-CH2-CH3}، البوتانون (ويسمى أيضًا
بوتان-2-ون). وللثاني حزمة \ce{O-H} ولا حزمة \ce{C=O}: فهو
\omterm{def:g11:functional-groups:alcohol}{كحول}، مثلًا بوتان-1-ول.
\end{solution}

\begin{solution}{exo:g11:infrared:7}
مجموعات \ce{O-H} في الحمض مرتبطة بروابط هيدروجينية قوية جدًا (إذ تتزاوج
\omterm{def:g7:atoms-and-molecules:molecule}{الجزيئات}، وترتبط كل \ce{O-H} بمجموعة \ce{C=O} في \omterm{def:g7:atoms-and-molecules:molecule}{الجزيء} الآخر)،
وكل \omterm{def:g7:atoms-and-molecules:molecule}{جزيء} مرتبط على نحو يختلف قليلًا: فتُرخى الرابطة
أكثر وتنتشر الحزمة على مجال واسع جدًا.
\end{solution}

\begin{solution}{exo:g11:infrared:8}
في السائل (A) حزمة \ce{O-H} عريضة، نحو
\qty{3350}{cm^{-1}}؛ وفي الغاز (B)، حيث \omterm{def:g7:atoms-and-molecules:molecule}{الجزيئات} متباعدة
ولا تكوّن روابط هيدروجينية، هي حزمة حادة نحو
\qty{3666}{cm^{-1}}.
\end{solution}

\begin{solution}{exo:g11:infrared:9}
بصعوبة: نحو \qty{1730}{cm^{-1}} \omterm{def:g11:functional-groups:carbonyl}{للألدهيد}
و\qty{1715}{cm^{-1}} \omterm{def:g11:functional-groups:carbonyl}{للكيتون}، وهما قيمتان متقاربتان. ومقارنة
الطيف كله، بما فيه \omterm{def:g11:infrared:band}{منطقة البصمة}، بأطياف مرجعية
هي التي تحسم، أو تقنية أخرى.
\end{solution}

\begin{solution}{exo:g11:infrared:10}
\omterm{def:g11:functional-groups:carboxylic-acid}{إستر} (حزمة \ce{C=O} نحو \qty{1735}{cm^{-1}}، وحزمة \ce{C-O} قوية، ولا
حزمة \ce{O-H})، مثلًا إيثانوات الإيثيل. وحمض الإيثانويك صيغته \ce{C2H4O2}، لا
\ce{C4H8O2}. ولحمض البوتانويك الصيغة الصحيحة لكنه يُظهر حزمة \ce{O-H}
العريضة جدًا الخاصة بالأحماض: فهو مستبعد.
\end{solution}

\begin{solution}{exo:g11:infrared:11}
بروبان-1-ول \omterm{def:g11:functional-groups:alcohol}{كحول}: فطيفه يشبه A. وحمض البروبانويك \omterm{def:g11:functional-groups:carboxylic-acid}{حمض كربوكسيلي}:
فطيفه يشبه D.
\end{solution}

\begin{solution}{exo:g11:infrared:12}
تتقلص حزمة \ce{O-H} العريضة نحو \qty{3350}{cm^{-1}} بينما تكبر حزمة
\ce{C=O} قوية نحو \qty{1715}{cm^{-1}}. ويكتمل التفاعل
حين تختفي حزمة \ce{O-H} وتكف حزمة \ce{C=O} عن
النمو.
\end{solution}

\begin{solution}{exo:g11:infrared:13}
يُظهر حمض البوتانويك حزمة \ce{O-H} عريضة جدًا من 2500 إلى
\qty{3300}{cm^{-1}}، ولا يُظهر \omterm{def:g11:functional-groups:carboxylic-acid}{الإستر} شيئًا منها؛ وتقع حزمة \ce{C=O} للحمض
نحو \qty{1710}{cm^{-1}}، أعرض، وحزمة \omterm{def:g11:functional-groups:carboxylic-acid}{الإستر} نحو \qty{1735}{cm^{-1}}،
حادة، مع حزمة \ce{C-O} قوية نحو \qty{1240}{cm^{-1}}. فللحمض
مجموعة \ce{O-H} مرتبطة بروابط هيدروجينية، ترخي أيضًا رابطته
\ce{C=O} قليلًا.
\end{solution}

\begin{solution}{exo:g11:infrared:14}
بقايا إيثانول (\omterm{def:g11:functional-groups:alcohol}{الكحول} الذي يُصنع منه \omterm{def:g11:functional-groups:carboxylic-acid}{الإستر}) أو ماء (من
\omterm{def:g4:air-a-mixture-of-gases:air}{الهواء} أو الغسل): ولكليهما مجموعات \ce{O-H} مرتبطة بروابط هيدروجينية. ويمكن
مقارنة العيّنة بطيف مرجعي، أو تجفيفها، أو تقطيرها، أو
قياس درجة غليانها.
\end{solution}

\begin{solution}{exo:g11:infrared:15}
\omterm{def:g10:lewis-and-shape:multiple-bond}{الرابطة الثنائية} نابض أشد صلابة من الرابطة الأحادية بين الذرتين
نفسيهما: فهي تهتز أسرع، عند \omterm{def:g11:infrared:wavenumber}{عدد موجي} أعلى (1715 مقابل
\qty{1050}{cm^{-1}}). وروابط \ce{O-H} و\ce{N-H} و\ce{C-H} كلها
تضم \omterm{def:g7:atoms-and-molecules:atom}{ذرة} هيدروجين، أخف \omterm{def:g7:atoms-and-molecules:atom}{الذرات} جميعًا: والكرات الخفيفة على نابض
تهتز بسرعة، ومن هنا الأعداد الموجية نحو 2900--\qty{3600}{cm^{-1}}.
\end{solution}

\begin{solution}{pb:g11:infrared:1}
\textbf{1.} \ce{CH3-CH2-OH}; \ce{CH3-CO-CH3}; \ce{CH3-COOH}.

\textbf{2.} \ce{C2H6O}؛ \ce{C3H6O}؛ \ce{C2H4O2}.

\textbf{3.} هيدروكسيل؛ كربونيل (\omterm{def:g11:functional-groups:carbonyl}{كيتون})؛ كربوكسيل.

\textbf{4.} الإيثانول: \ce{O-H} و\ce{C-H}. البروبانون: \ce{C-H}
و\ce{C=O}. حمض الإيثانويك: \ce{O-H} و\ce{C-H} و\ce{C=O}.

\textbf{5.} لا حزمة \ce{C=O}، وحزمة \ce{O-H} عريضة: الإيثانول.

\textbf{6.} حزمة \ce{O-H} عريضة جدًا من 2500 إلى \qty{3300}{cm^{-1}}
وحزمة \ce{C=O} قوية نحو \qty{1710}{cm^{-1}}: حمض الإيثانويك.

\textbf{7.} البروبانون: حزمة \ce{C=O} القوية جدًا عند
\qty{1715}{cm^{-1}}، دون حزمة \ce{O-H}.

\textbf{8.} على الأرجح (خل، \omterm{def:g6:solutions-and-solubility:solute}{مذيب}، \omterm{def:g11:functional-groups:alcohol}{كحول})، لكن شمّ سوائل مجهولة
خطر؛ والطيف آمن وأكيد.

\textbf{9.} مجموعات \ce{O-H} في الحمض مرتبطة بروابط هيدروجينية أقوى
(\omterm{def:g7:atoms-and-molecules:molecule}{جزيئات} متزاوجة)، فتكون الحزمة أدنى وأعرض بكثير.

\textbf{10.} \ce{CH3-O-CH3}، \ce{C2H6O}؛ ليس فيه \ce{O-H}، فلا حزمة
\ce{O-H}.

\textbf{11.} حزمة \ce{C=O}: نحو \qty{1730}{cm^{-1}} للبروبانال،
ونحو \qty{1715}{cm^{-1}} للبروبانون.

\textbf{12.} حزمة \ce{O-H}: فليس في \omterm{def:g11:functional-groups:carboxylic-acid}{الإستر} \ce{O-H}. وتقع حزمته \ce{C=O}
نحو \qty{1735}{cm^{-1}}.

\textbf{13.} \omterm{def:g11:organic-skeletons:isomer}{للمتماكبات} صيغة واحدة؛ لكن مجموعاتها الوظيفية تختلف،
والطيف يُظهر المجموعات.

\textbf{14.} \qty{1715}{cm^{-1}}.

\textbf{15.} $1715 \times 100 = \qty{1.715e5}{m^{-1}}$.

\textbf{16.} $\lambda = 1/\num{1.715e5} = \qty{5.83e-6}{m} =
\qty{5.83}{\micro m}$.

\textbf{17.} لا: يمتد الضوء المرئي من نحو 0.4 إلى
\qty{0.75}{\micro m}. إنه إشعاع تحت أحمر.

\textbf{18.} $\num{1.05e5}\ \si{m^{-1}}$، إذن $\lambda = \qty{9.52e-6}{m}
= \qty{9.52}{\micro m}$.

\textbf{19.} الرابطة \ce{C=O} (\omterm{def:g11:infrared:wavenumber}{عدد موجي} أعلى، وطول موجي أقصر).

\textbf{20.} نحو \qty{5.8}{\micro m}.
\end{solution}
